\usepackage[all]{xy}
\usepackage{bm}
\usepackage{dcolumn}
\usepackage{tikz-cd}

 \newcolumntype{2}{D{.}{}{2.0}}
 \xyoption{2cell}

\newcommand{\hiddenpower}[2] { \ifnum \numexpr#2=1 #1 \else #1^#2 \fi }
\numberwithin{equation}{section}

\def\be{\begin{equation}}
\def\ee{\end{equation}}
\def\ba{\begin{eqnarray}}
\def\ea{\end{eqnarray}}
\def\db{\bar{\partial}}
\def\Jb{\bar{J}}
\def\non{\nonumber}
\def\del{\partial}

\newcommand{\dd}{\text{d}}
\newcommand{\pd}{\partial}

\usepackage{xparse}
\ExplSyntaxOn
\newcounter{diff_order}
\newcounter{diff_power}

\newcommand{\rawdiff}[3]
{
	% Get the order of the differential.
	\setcounter{diff_order}{0}
	\clist_map_inline:nn{#3}{\stepcounter{diff_order}}
	
	\frac{\hiddenpower{#1}{\thediff_order} #2}
	{
		\def\old_var{DefaultValue}
		\setcounter{diff_power}{0}
		
		% Iterate through the comma separated list in #3.
		\clist_map_inline:nn{#3}
		{
			\def\new_var{##1}
			% Skip the first one (as it's currently just DefaultValue).
			\ifnum \thediff_power=0
				\stepcounter{diff_power}
			\else
				% If you're the same as last time, count up. Else, print.
				\tl_if_eq:NNTF \new_var \old_var
				{\stepcounter{diff_power}}
				{
					#1 \hiddenpower{\old_var}{\thediff_power}
					\setcounter{diff_power}{1}
				}
			\fi

			\def\old_var{##1}
		}
		
		% Finally, print what's left.
		#1 \hiddenpower{\old_var}{\thediff_power}
	}
}

\def\Label#1{\label{#1}\ifmmode\llap{[#1] }\else
 \marginpar{\smash{\hbox{\tiny [#1]}}}\fi}
 \def\Label{\label}
% Totally worth it!
\newcommand{\diff}[2]{\rawdiff{\dd}{#1}{#2}}
\newcommand{\pdiff}[2]{\rawdiff{\pd}{#1}{#2}}

\ExplSyntaxOff

\newcommand{\lb}{\left(}
\newcommand{\rb}{\right)}

\newtheorem{thm}{Theorem}[section]
\newtheorem{lem}[thm]{Lemma}
\newtheorem{cor}[thm]{Corollary}
\newtheorem{pro}[thm]{Proposition}
\newtheorem{conj}[thm]{Conjecture}

\theoremstyle{definition}
\newtheorem{defn}[thm]{Definition}
\newtheorem{problem}[thm]{Problem}
\newtheorem{notation}[thm]{Notation}
\newtheorem{question}[thm]{Question}
\newtheorem{convention}[thm]{Convention}
\newtheorem{rem}[thm]{Remark}
\newtheorem{exam}[thm]{Example}

\newcommand{\Aut}{\operatorname{Aut}}
\newcommand{\Sym}{\operatorname{Sym}}
\newcommand{\Z}{\mathbb{Z}}
\newcommand{\id}{\operatorname{id}}

\DeclareMathOperator{\EEnd}{End}%\newenvironment{proof}{\par\noindent{\bf Proof.}}{$\qed$\par\bigskip}
%\newcommand{\qed}{\enspace\vrule height6pt width4pt depth2pt}

\renewcommand{\sin}[2][1]{\hiddenpower{\text{sin}}{#1} \lb #2 \rb}
\newcommand{\asin}[2][1]{\hiddenpower{\text{asin}}{#1} \lb #2 \rb}
\renewcommand{\cos}[2][1]{\hiddenpower{\text{cos}}{#1} \lb #2 \rb}
\newcommand{\acos}[2][1]{\hiddenpower{\text{acos}}{#1} \lb #2 \rb}
\renewcommand{\tan}[2][1]{\hiddenpower{\text{tan}}{#1} \lb #2 \rb}
\newcommand{\atan}[2][1]{\hiddenpower{\text{atan}}{#1} \lb #2 \rb}
\renewcommand{\csc}[2][1]{\hiddenpower{\text{csc}}{#1} \lb #2 \rb}
\renewcommand{\sec}[2][1]{\hiddenpower{\text{sec}}{#1} \lb #2 \rb}
\renewcommand{\cot}[2][1]{\hiddenpower{\text{cot}}{#1} \lb #2 \rb}
\renewcommand{\sinh}[2][1]{\hiddenpower{\text{sinh}}{#1} \lb #2 \rb}
\renewcommand{\cosh}[2][1]{\hiddenpower{\text{cosh}}{#1} \lb #2 \rb}
\renewcommand{\tanh}[2][1]{\hiddenpower{\text{tanh}}{#1} \lb #2 \rb}
\newcommand{\csch}[2][1]{\hiddenpower{\text{csch}}{#1} \lb #2 \rb}
\newcommand{\sech}[2][1]{\hiddenpower{\text{sech}}{#1} \lb #2 \rb}
\renewcommand{\coth}[2][1]{\hiddenpower{\text{coth}}{#1} \lb #2 \rb}
\renewcommand{\ln}[1]{\text{ln} \lb #1 \rb}
\newcommand{\e}[1]{\text{e}^{#1}}

\newcommand{\tr}[1]{\text{tr}\left\lbrace #1 \right\rbrace}
\newcommand{\ptr}[2]{\text{tr}_{#1}\left\lbrace #2 \right\rbrace}

\newcommand{\rp}[1]{\ensuremath{\mathbb{R} P^{#1}}}
\newcommand{\cp}[1]{\ensuremath{\mathbb{C} P^{#1}}}
\newcommand{\set}[1]{\ensuremath{\left\lbrace #1 \right\rbrace}}
\newcommand{\ptset}{\ensuremath{\lbrace \text{pt.} \rbrace}}
\newcommand{\gen}[1]{\ensuremath{\left\langle #1 \right\rangle}}

\newcommand{\comm}[2]{\ensuremath{\left[ #1, #2 \right]}}
\newcommand{\pcomm}[2]{\ensuremath{\set{#1, #2}}}

\newcommand{\nbym}[2]{\ensuremath{#1{\times}#2}}
\newcommand{\nbyn}[1]{\nbym{#1}{#1}}

\newcommand{\gln}[1]{\ensuremath{\mathfrak{gl}_{#1}}}
\newcommand{\sln}[1]{\ensuremath{\mathfrak{sl}_{#1}}}

% Preserve original author display pagination from after document start.
\allowdisplaybreaks[4]

% PDF outline compatibility only; printed author headings remain unchanged.
\pdfstringdefDisableCommands{%
  \def\_{ }%
  \def\boldsymbol#1{#1}%
  \def\mathfrak#1{#1}%
  \def\mathcal#1{#1}%
}
\hypersetup{bookmarks=false}
